\documentclass[11pt]{article}
\usepackage[margin=1in]{geometry}
\usepackage{amsmath,amssymb,amsthm}
\usepackage[T1]{fontenc}
\usepackage{lmodern}
\usepackage[hidelinks]{hyperref}

\newcommand{\Lip}{\operatorname{Lip}_0}

\title{Literature Status: Local Lipschitz Functions Are Daugavet Points}
\author{fa\_banach\_001 lane 0}
\date{2026-07-18}

\begin{document}
\maketitle

\section*{Status}

\textbf{Classification.} \texttt{literature\_already\_answered}.

\textbf{Source paper.} Mingu Jung and Abraham Rueda Zoca,
\emph{Daugavet points and $\Delta$-points in Lipschitz-free spaces},
arXiv:2010.09357.

\textbf{Supporting paper.} Rainis Haller, Andre Ostrak, and Mart Poldvere,
\emph{Diameter two properties for spaces of Lipschitz functions},
arXiv:2205.13287.

\section*{Source Question}

In Section 3 of arXiv:2010.09357, after Proposition 3.4, Jung and Rueda Zoca
prove that spreadingly local functions in $S_{\Lip(M)}$ are Daugavet points.
They then ask whether the same conclusion remains true under the weaker
assumption that $f$ is merely local. The same section proves the weaker
weak-star version: local $f\in S_{\Lip(M)}$ are $w^*$-Daugavet points.

\section*{Supporting Answer}

The later paper arXiv:2205.13287 explicitly identifies this Jung--Rueda Zoca
question. In its introduction, Theorem 1.5 states:
\[
  \text{if } M \text{ is a pointed metric space and } f\in S_{\Lip(M)}
  \text{ is local, then } f \text{ is a Daugavet point.}
\]
Section 5, titled ``Local norm-one Lipschitz function is a Daugavet point'',
repeats the source question and says that the theorem fully answers it.

\section*{Mechanism of the Answer}

The decisive improvement is Proposition 5.3 of arXiv:2205.13287. For every
$F\in S_{\Lip(M)^*}$ and every sequence of distinct pairs $(u_n,v_n)$ with
$d(u_n,v_n)\to0$, it proves
\[
  \|F+m_{u_n,v_n}\|\longrightarrow 2 .
\]
This is the full weak-slice analogue of the free-space estimate used in the
source paper's weak-star proof. The proof represents arbitrary weak
functionals on $\Lip(M)$ by bounded finitely additive measures on the de
Leeuw transform
\[
  \widetilde M=(M\times M)\setminus\{(x,x):x\in M\}.
\]
It then modifies a test Lipschitz function on a small metric set whose first
and second coordinate cylinders have arbitrarily small measure. This preserves
membership in an arbitrary weak slice while forcing almost maximal opposite
slope on one of the small pairs supplied by locality.

\section*{Conclusion}

The exact question extracted from arXiv:2010.09357 has an affirmative later
literature answer:
\[
  \boxed{\text{Every local norm-one } f\in\Lip(M) \text{ is a Daugavet point.}}
\]
This packet is therefore provenance/status memory rather than a new result of
the run.

\section*{Search and Novelty Notes}

The identification was found by searching the run indexes and parsed arXiv
sources for \texttt{2010.09357}, \texttt{local Lipschitz},
\texttt{spreadingly local}, \texttt{w* Daugavet}, and \texttt{Daugavet point}.
The decisive local source hit is arXiv:2205.13287, which explicitly cites
Jung--Rueda Zoca and says the theorem fully answers the question. This places
the packet in \texttt{solutions/literature\_already\_answered/}, not in
\texttt{solutions/full/}.

\end{document}
